\input{../tex-inputs/latex-leseansicht-vorspann}

\section[Arthur Schnitzler an Gustav Schwarzkopf, 3. 7. 1923]{L04499 Arthur Schnitzler an Gustav Schwarzkopf, 3. 7. 1923}
\nopagebreak\mylabel{L04499v}
\rehead{ }\normalsize\beginnumbering\briefempfaengerindex{Schwarzkopf, Gustav@\textsc{Schwarzkopf, Gustav}!zzzSchnitzler, Arthur@\emph{von Arthur Schnitzler}!1923-07-031@{3. 7. 1923}|(be}
\toendnotes[C]{\smallbreak\pagebreak[2]}
\correspDesc{Versand  durch Arthur Schnitzler am 3. 7. 1923 in Baden-Baden
\newline{}Übermittlung  am 3. 7. 1923 in Baden-Baden
\newline{}Erhalt  durch Gustav Schwarzkopf im Zeitraum [4. 7. 1923 – 8. 7. 1923?] in Wien}\toendnotes[C]{\smallbreak}
\Standort{DLA, A:Schnitzler, HS.1985.1.1897.}
\physDesc{Bildpostkarte, 137 Zeichen
\newline{}Handschrift: Bleistift, deutsche Kurrent
\newline{}Versand: 1) Stempel: »\nobreak{}\oindex{Baden-Baden@\textbf{Baden-Baden}|pwk}Baden-Baden, 3. 7. 23, 4–5 N\nobreak{}«.   2) Stempel: »\nobreak{}Baden-Badener Rennen\orgindex{Grosser Preis von Baden [Pferderennen]@Grosser Preis von Baden [Pferderennen]|pw} / 24. 26. 28. 31. Aug. / \textcolor{gray}{02.} Sept\nobreak{}«. }\toendnotes[C]{\smallbreak}
\pstart
           \noindent{}{\pb}\textcolor{gray}{\textbf{BADEN-BADEN\oindex{Baden-Baden@\textbf{Baden-Baden}|pw}, von Pension Hohenstein\oindex{Villa Hohenstein@\textbf{Villa Hohenstein}|pw} gesehen}}\pend
           \pstart{}{\pb}\textsc{Herrn}\pend{}\pstart{}Guſtav \textsc{Schwarzkopf}\pend{}\pstart{}\textsc{Wien I}\oindex{I., Innere Stadt@\textbf{I., Innere Stadt}|pw}\pend{}\pstart{}\textsc{Tiefer Graben 17}\oindex{Wien@\textbf{Wien}!I., Innere Stadt@\textbf{I., Innere Stadt}!Tiefer Graben 17@\textbf{Tiefer Graben 17}|pw}\pend{}{\bigskip}\vspace{1em}
\pstart
           \raggedleft{}Baden-Baden\oindex{Baden-Baden@\textbf{Baden-Baden}|pw}{ }3. 7. 23\pend
           \vspace{0.5em}
\pstart
           Herzliche Grüße, auch
      für Ihren Bruder\pwindex{Schwarzkopf, Max 12.\,6.\,1857 Wien – 14.\,4.\,1928 ebd.@\textsc{Schwarzkopf, Max} (12.\,6.\,1857 Wien – 14.\,4.\,1928 ebd.)|pwv}! u
      auf baldgs Wiederſehen!\pend
           \pstart Ihr \spacefill\mbox{Arthur}\pend{}\selectlanguage{ngerman}\endnumbering\briefempfaengerindex{Schwarzkopf, Gustav@\textsc{Schwarzkopf, Gustav}!zzzSchnitzler, Arthur@\emph{von Arthur Schnitzler}!1923-07-031@{3. 7. 1923}|)be}\mylabel{L04499h}
\begin{anhang}
\end{anhang}\newcommand{\dateiname}{L04499}\newcommand{\titel}{Arthur Schnitzler an Gustav Schwarzkopf, 3. 7. 1923}\newcommand{\editorInnen}{Herausgegeben von Jahnke, SelmaMüller, Martin Anton}\input{../tex-inputs/latex-leseansicht-abspann}
